% TOP_PROBLEM: {"id": 30001464, "title": "Rigidity and Assembly Conjectures for Discrete Groups", "category_id": 17, "category": {"id": 17, "name": "group_theory", "display_name": "Group Theory", "description": "Problems about groups, group actions, representations, and related algebraic structures.", "slug": "group-theory", "order_index": 17, "created_at": "2026-07-31T15:26:25.671Z"}, "status": "partially_solved"}
% FIRST_POSTED: 2026-10-06
% ATTEMPT_STATUS: unresolved
% OPENAI_CLAIM: {"schema_version":1,"source_repo":"https://github.com/openai/math","source_commit":"adc7f1241b42e322a6451854ab7e4b4c146bf78a","release_date":"2026-10-06","match":"partial","resolution":"partial","independently_reviewed":false,"second_pass":"n/a","families":[{"family":"207","title":"The ℓ¹-Bass conjecture for all discrete groups","subject":"Algebra","manuscripts":[{"title":"The Bass trace conjecture and the characteristic-zero Kaplansky idempotent conjecture","pdf_path":"preprints/The-Bass-trace-conjecture-for-complex-group-rings-September-24-2026/The-Bass-trace-conjecture-for-complex-group-rings-September-24-2026.pdf","readme_path":"preprints/The-Bass-trace-conjecture-for-complex-group-rings-September-24-2026/README.md","date":"2026-09-24","theorem_ref":"thm:bass (sections/01-introduction.tex:45); cor:torsion-free (sections/01-introduction.tex:98)","statement_sources":[{"path":"preprints/The-Bass-trace-conjecture-for-complex-group-rings-September-24-2026/build/sections/01-introduction.tex","line":45,"label":"thm:bass"},{"path":"preprints/The-Bass-trace-conjecture-for-complex-group-rings-September-24-2026/build/sections/01-introduction.tex","line":98,"label":"cor:torsion-free"}]}],"lean":{"doc_path":"lean/docs/207.md","comparators":["lean/ComparatorChallenges/BassTorsionFree.json","lean/ComparatorChallenges/BassTorsionFree.lean","lean/ComparatorChallenges/BassTrace.json","lean/ComparatorChallenges/BassTrace.lean"],"declarations":["OAI.BassTrace.RightProjective.bassTraceModules_vanishing_and_support"],"covers_main_theorem":true},"reasoning_trace":null},{"family":"246","title":"Cannon's conjecture","subject":"Group theory","manuscripts":[{"title":"A Modulus Proof of Cannon's Conjecture","pdf_path":"preprints/A-Modulus-Proof-of-Cannons-Conjecture-September-23-2026/paper.pdf","readme_path":"preprints/A-Modulus-Proof-of-Cannons-Conjecture-September-23-2026/README.md","date":"2026-09-23","theorem_ref":"thm:main (main.tex:94); cor:torsion-free (main.tex:277); cor:virtual-structure (main.tex:289)","statement_sources":[{"path":"preprints/A-Modulus-Proof-of-Cannons-Conjecture-September-23-2026/build/main.tex","line":94,"label":"thm:main"},{"path":"preprints/A-Modulus-Proof-of-Cannons-Conjecture-September-23-2026/build/main.tex","line":277,"label":"cor:torsion-free"},{"path":"preprints/A-Modulus-Proof-of-Cannons-Conjecture-September-23-2026/build/main.tex","line":289,"label":"cor:virtual-structure"}]}],"lean":{"doc_path":"lean/docs/246.md","comparators":["lean/ComparatorChallenges/CannonGeometricAction.json","lean/ComparatorChallenges/CannonGeometricAction.lean"],"declarations":[],"covers_main_theorem":true},"reasoning_trace":null},{"family":"320","title":"Nonhomeomorphic closed aspherical four-manifolds","subject":"Topology","manuscripts":[{"title":"Nonhomeomorphic closed aspherical four-manifolds with the same homotopy type","pdf_path":"preprints/Nonhomeomorphic-closed-aspherical-four-manifolds-with-the-same-homotopy-type-October-4-2026/paper.pdf","readme_path":"preprints/Nonhomeomorphic-closed-aspherical-four-manifolds-with-the-same-homotopy-type-October-4-2026/README.md","date":"2026-10-04","theorem_ref":"thm:pair (sections/introduction.tex:12)","statement_sources":[{"path":"preprints/Nonhomeomorphic-closed-aspherical-four-manifolds-with-the-same-homotopy-type-October-4-2026/build/sections/introduction.tex","line":12,"label":"thm:pair"}]}],"lean":null,"reasoning_trace":null}]}
\documentclass[11pt]{article}
\usepackage[margin=25mm]{geometry}
\usepackage{fontspec}
\setmainfont{Latin Modern Roman}
\usepackage{amsmath,amssymb,mathtools}
\usepackage{unicode-math}
\setmathfont{Latin Modern Math}
\usepackage{enumitem}
\usepackage{xurl,hyperref,bookmark}
\hypersetup{colorlinks=true,urlcolor=blue}
\setcounter{secnumdepth}{0}
\setlength{\parindent}{0pt}
\setlength{\parskip}{5pt}
\setlength{\emergencystretch}{3em}
\begin{document}
\section*{Rigidity and Assembly Conjectures for Discrete Groups}

\subsection{OpenAI's claim}
The linked manuscripts present the following claims by OpenAI: The Bass trace restriction is claimed for complex group rings of arbitrary discrete groups, implying triviality of scalar idempotents for torsion-free groups over characteristic-zero domains; A hyperbolic group with two-sphere boundary is claimed to act geometrically on hyperbolic three-space with finite kernel; Two homotopy-equivalent closed aspherical four-manifolds with hyperbolic fundamental group are claimed nonhomeomorphic, with an additional nongeometric self-homotopy equivalence. Together they provide the following support within the scopes described below.

\subsection{Mathematical statement of the claimed result}
\paragraph{Result 1.} Let \(G\) be word-hyperbolic with Gromov boundary \(\partial G\cong S^2\). The source claims a homomorphism
\[
 \rho:G\longrightarrow\operatorname{Isom}(\mathbb H^3)
\]
whose action on hyperbolic three-space is proper and cocompact, and whose kernel is finite. With the additional assumption that \(G\) is torsion-free, the corollaries identify \(G\cong\pi_1(M)\), where \(M\) is closed and has a hyperbolic metric in dimension three, and supply a finite-index subgroup
\[
 H\cong\pi_1(\Sigma)\rtimes\mathbb Z,
\]
where \(\Sigma\) is a closed connected orientable surface. They also assert residual finiteness, virtual compact specialness, and positive virtual first Betti number in that torsion-free case.

\paragraph{Result 2.} The construction claims two closed connected aspherical topological four-manifolds \(M,N\) with
\[
 M\simeq N,\qquad M\not\cong N\text{ by homeomorphism},\qquad
 \pi_1M\cong\pi_1N\text{ word-hyperbolic}.
\]
Thus homotopy equivalence and hyperbolicity do not force topological homeomorphism for this claimed four-dimensional pair.

\paragraph{Result 3.} For an arbitrary discrete group \(G\), let \(\operatorname{HS}_G\) be the Hattori--Stallings trace, whose value at a conjugacy class sums the coefficients on that class. For every \(x\in K_0(\mathbb C[G])\),
\[
 |g|=\infty\quad\Longrightarrow\quad\operatorname{HS}_G(x)([g])=0.
\]
The torsion-free corollary states that, for every \(q\ge1\) and idempotent \(e\in M_q(\mathbb C[G])\), its identity-coefficient trace equals the rank of its augmentation matrix:
\[
 \tau_{G,q}(e)=\operatorname{rank}_{\mathbb C}(\epsilon_{G,q}(e)).
\]
For such \(G\), any commutative unital domain \(R\) of characteristic zero satisfies
\[
 e\in R[G],\ e^2=e\quad\Longrightarrow\quad e\in\{0,1\}.
\]
The characteristic-zero qualification does not cover positive-characteristic group rings.

\subsection{Outline of the argument}
An idempotent is encoded by finite cyclic path configurations whose weighted deletion relations generate cycles. Exterior differential forms and a cyclic Pfaffian identity detect forbidden conjugacy support; trace-rank arguments then turn the complex trace conclusion into the stated idempotent consequence.

The argument applies a boundary modulus criterion. Unbounded modulus produces scalar limits analyzed by exponential records or subexponential probability flows; comparing same-depth expansions forces the required measure control and the three-dimensional hyperbolic action.

Reflected cyclic covers realize the relevant symmetries while preserving asphericity. A parity-chamber construction and cyclic-cover criteria distinguish manifold models despite the common homotopy type, and the realized symmetry gives the self-equivalence obstruction.

\subsection{Scope relative to this problem}
Corollary1.2(ii) proves the idempotent assertion for characteristic-zero domains, but leaves positive characteristic and the independent rigidity and assembly assertions unanswered. The main theorem and its closed-hyperbolic-three-manifold consequences settle the S\textasciicircum{}2-boundary subcase. The record quantifies over additional boundary types, general hyperbolic groups, dimensions, or independent bundle clauses that this paper does not settle. The nonhomeomorphic aspherical pair refutes the Borel-rigidity clause of this broad bundle, but does not decide its idempotent, K-group, signature, FP/FF or higher-dimensional realization clauses. This entry records partial support; complete coverage of the question is not established by this catalogue.

\subsection{Formal verification}
For family 207, the documented Lean statements give complex group-ring Bass trace support and characteristic-zero torsion-free trace rank and idempotent triviality. The separate $\ell^1$ theorem is not selected. The listed Comparator files are \texttt{BassTorsionFree.lean}, \texttt{BassTrace.lean}. Named declarations include \texttt{OAI.BassTrace.RightProjective.bassTraceModules\_vanishing\_and\_support}. For family 246, the documented theorem gives a proper cocompact hyperbolic-three-space action with finite kernel from the two-sphere boundary hypothesis. The listed Comparator files are \texttt{CannonGeometricAction.lean}. The inventory does not supply a named declaration for these Comparator entries. Family 320 has no Lean scope page in this release. Its selected result is not presented as formalized. This site has not executed the Lean code. The scope descriptions do not certify the remaining manuscript arguments.

\subsection{Status caveat}
OpenAI attributes these manuscripts to work by an internal model. Its release cautions that arguments without formalization can contain mistakes. This catalogue has not independently checked the proofs; the labels report the mathematical claims attributed to OpenAI.

\subsection{Sources}
\begin{itemize}
\item OpenAI, \emph{The Bass trace conjecture and the characteristic-zero Kaplansky idempotent conjecture}, 2026-09-24; citation key \texttt{OAI:The-Bass-trace-conjecture-for-complex-group-rings-September-24-2026}. \href{https://github.com/openai/math/blob/adc7f1241b42e322a6451854ab7e4b4c146bf78a/preprints/The-Bass-trace-conjecture-for-complex-group-rings-September-24-2026/The-Bass-trace-conjecture-for-complex-group-rings-September-24-2026.pdf}{PDF (GitHub)}; \href{https://github.com/openai/math/raw/adc7f1241b42e322a6451854ab7e4b4c146bf78a/preprints/The-Bass-trace-conjecture-for-complex-group-rings-September-24-2026/The-Bass-trace-conjecture-for-complex-group-rings-September-24-2026.pdf}{PDF (direct)}; \href{https://github.com/openai/math/blob/adc7f1241b42e322a6451854ab7e4b4c146bf78a/preprints/The-Bass-trace-conjecture-for-complex-group-rings-September-24-2026/README.md}{Manuscript page and citation}; \href{https://github.com/openai/math/blob/main/preprints/The-Bass-trace-conjecture-for-complex-group-rings-September-24-2026/The-Bass-trace-conjecture-for-complex-group-rings-September-24-2026.pdf}{Latest version}; Statement source: \href{https://github.com/openai/math/blob/adc7f1241b42e322a6451854ab7e4b4c146bf78a/preprints/The-Bass-trace-conjecture-for-complex-group-rings-September-24-2026/build/sections/01-introduction.tex\#L45}{thm:bass}; Statement source: \href{https://github.com/openai/math/blob/adc7f1241b42e322a6451854ab7e4b4c146bf78a/preprints/The-Bass-trace-conjecture-for-complex-group-rings-September-24-2026/build/sections/01-introduction.tex\#L98}{cor:torsion-free}.
\item Family 207: \href{https://github.com/openai/math/blob/adc7f1241b42e322a6451854ab7e4b4c146bf78a/lean/docs/207.md}{Lean scope}.
\item \href{https://github.com/openai/math/blob/adc7f1241b42e322a6451854ab7e4b4c146bf78a/lean/ComparatorChallenges/BassTorsionFree.json}{Comparator BassTorsionFree.json}.
\item \href{https://github.com/openai/math/blob/adc7f1241b42e322a6451854ab7e4b4c146bf78a/lean/ComparatorChallenges/BassTorsionFree.lean}{Comparator BassTorsionFree.lean}.
\item \href{https://github.com/openai/math/blob/adc7f1241b42e322a6451854ab7e4b4c146bf78a/lean/ComparatorChallenges/BassTrace.json}{Comparator BassTrace.json}.
\item \href{https://github.com/openai/math/blob/adc7f1241b42e322a6451854ab7e4b4c146bf78a/lean/ComparatorChallenges/BassTrace.lean}{Comparator BassTrace.lean}.
\item OpenAI, \emph{A Modulus Proof of Cannon's Conjecture}, 2026-09-23; citation key \texttt{OAI:A-Modulus-Proof-of-Cannons-Conjecture-September-23-2026}. \href{https://github.com/openai/math/blob/adc7f1241b42e322a6451854ab7e4b4c146bf78a/preprints/A-Modulus-Proof-of-Cannons-Conjecture-September-23-2026/paper.pdf}{PDF (GitHub)}; \href{https://github.com/openai/math/raw/adc7f1241b42e322a6451854ab7e4b4c146bf78a/preprints/A-Modulus-Proof-of-Cannons-Conjecture-September-23-2026/paper.pdf}{PDF (direct)}; \href{https://github.com/openai/math/blob/adc7f1241b42e322a6451854ab7e4b4c146bf78a/preprints/A-Modulus-Proof-of-Cannons-Conjecture-September-23-2026/README.md}{Manuscript page and citation}; \href{https://github.com/openai/math/blob/main/preprints/A-Modulus-Proof-of-Cannons-Conjecture-September-23-2026/paper.pdf}{Latest version}; Statement source: \href{https://github.com/openai/math/blob/adc7f1241b42e322a6451854ab7e4b4c146bf78a/preprints/A-Modulus-Proof-of-Cannons-Conjecture-September-23-2026/build/main.tex\#L94}{thm:main}; Statement source: \href{https://github.com/openai/math/blob/adc7f1241b42e322a6451854ab7e4b4c146bf78a/preprints/A-Modulus-Proof-of-Cannons-Conjecture-September-23-2026/build/main.tex\#L277}{cor:torsion-free}; Statement source: \href{https://github.com/openai/math/blob/adc7f1241b42e322a6451854ab7e4b4c146bf78a/preprints/A-Modulus-Proof-of-Cannons-Conjecture-September-23-2026/build/main.tex\#L289}{cor:virtual-structure}.
\item Family 246: \href{https://github.com/openai/math/blob/adc7f1241b42e322a6451854ab7e4b4c146bf78a/lean/docs/246.md}{Lean scope}.
\item \href{https://github.com/openai/math/blob/adc7f1241b42e322a6451854ab7e4b4c146bf78a/lean/ComparatorChallenges/CannonGeometricAction.json}{Comparator CannonGeometricAction.json}.
\item \href{https://github.com/openai/math/blob/adc7f1241b42e322a6451854ab7e4b4c146bf78a/lean/ComparatorChallenges/CannonGeometricAction.lean}{Comparator CannonGeometricAction.lean}.
\item OpenAI, \emph{Nonhomeomorphic closed aspherical four-manifolds with the same homotopy type}, 2026-10-04; citation key \texttt{OAI:Nonhomeomorphic-closed-aspherical-four-manifolds-with-the-same-homotopy-type-October-4-2026}. \href{https://github.com/openai/math/blob/adc7f1241b42e322a6451854ab7e4b4c146bf78a/preprints/Nonhomeomorphic-closed-aspherical-four-manifolds-with-the-same-homotopy-type-October-4-2026/paper.pdf}{PDF (GitHub)}; \href{https://github.com/openai/math/raw/adc7f1241b42e322a6451854ab7e4b4c146bf78a/preprints/Nonhomeomorphic-closed-aspherical-four-manifolds-with-the-same-homotopy-type-October-4-2026/paper.pdf}{PDF (direct)}; \href{https://github.com/openai/math/blob/adc7f1241b42e322a6451854ab7e4b4c146bf78a/preprints/Nonhomeomorphic-closed-aspherical-four-manifolds-with-the-same-homotopy-type-October-4-2026/README.md}{Manuscript page and citation}; \href{https://github.com/openai/math/blob/main/preprints/Nonhomeomorphic-closed-aspherical-four-manifolds-with-the-same-homotopy-type-October-4-2026/paper.pdf}{Latest version}; Statement source: \href{https://github.com/openai/math/blob/adc7f1241b42e322a6451854ab7e4b4c146bf78a/preprints/Nonhomeomorphic-closed-aspherical-four-manifolds-with-the-same-homotopy-type-October-4-2026/build/sections/introduction.tex\#L12}{thm:pair}.
\item \href{https://github.com/openai/math/blob/adc7f1241b42e322a6451854ab7e4b4c146bf78a/CONTENTS.md}{OpenAI catalogue entry}.
\end{itemize}
\end{document}
